\documentclass[12pt,a4paper]{article}

\usepackage[utf8]{inputenc}
\usepackage[T1]{fontenc}
\usepackage[brazil]{babel}
\usepackage[margin=2.5cm]{geometry}
\usepackage{graphicx}
\usepackage{listings}
\usepackage{xcolor}
\usepackage{hyperref}
\usepackage{amsmath,amssymb}
\usepackage{tikz}
\usetikzlibrary{positioning,arrows.meta,shapes,shadows}
\usepackage{booktabs}
\usepackage{caption}
\usepackage{subcaption}
\usepackage{fancyhdr}
\usepackage{setspace}
\usepackage{tocloft}

% ============================================================================
% CONFIGURAÇÕES DE CÓDIGO
% ============================================================================
\lstset{
  basicstyle=\ttfamily\small,
  keywordstyle=\color{blue!70!black}\bfseries,
  commentstyle=\color{gray},
  stringstyle=\color{green!50!black},
  showstringspaces=false,
  breaklines=true,
  numbers=left,
  numberstyle=\tiny\color{gray},
  frame=single,
  language=Python,
  tabsize=2,
  showspaces=false,
  showtabs=false,
  breakatwhitespace=true
}

% ============================================================================
% ESPAÇAMENTO
% ============================================================================
\setstretch{1.5}
\setlength{\parindent}{1.27cm}

% ============================================================================
% AJUSTES DE LAYOUT
% ============================================================================
\hbadness=10000  % Suprimir avisos de caixas muito largas (overfull/underfull hbox)
\vbadness=10000  % Suprimir avisos verticais

% Aumentar tolerância para quebras de linha
\hyphenpenalty=10000
\exhyphenpenalty=10000

% Ajustar espaçamento do sumário
\setlength{\cftchapnumwidth}{2em}
\setlength{\cftsecnumwidth}{2.5em}
\setlength{\cftsubsecnumwidth}{3em}

% Permitir quebras de linha no sumário
\let\original@cftdotfill\cftdotfill
\renewcommand{\cftdotfill}[1]{\cfttabfillnum{#1}}

% ============================================================================
% CABEÇALHO E RODAPÉ
% ============================================================================
\pagestyle{fancy}
\fancyhf{}
\fancyhead[R]{\small \textit{Análise de Colaboração em Repositórios GitHub}}
\fancyfoot[C]{\thepage}
\renewcommand{\headrulewidth}{0.4pt}

% ============================================================================
% TÍTULO
% ============================================================================
\title{%
  \textbf{\Large Análise da Rede de Colaboração em Repositórios do GitHub}\\[0.5cm]
  \textbf{\large utilizando Teoria de Grafos}\\[1cm]
  \normalsize Trabalho Prático da Disciplina de Teoria de Grafos e Computabilidade\\[0.3cm]
  \normalsize Pontifícia Universidade Católica de Minas Gerais\\[0.3cm]
  \normalsize Curso de Engenharia de Software\\[0.3cm]
  \normalsize Professor: Leonardo V. Cardoso\\[0.5cm]
}

\author{%
  \textbf{Participantes do Projeto:}\\[0.3cm]
  Davi Nunes Carvalho\\[0.2cm]
  Luiz Fernando Batista Moreira\\[0.2cm]
  Josué Carlos Goulart dos Reis\\[0.2cm]
  João Victor Russo Marquito%
}

\date{%
  \vspace{0.5cm}
  \textbf{2026/1}\\[0.3cm]
  Junho de 2026%
}

\begin{document}

\maketitle

% ============================================================================
% RESUMO
% ============================================================================
\begin{abstract}
  Este trabalho apresenta uma ferramenta completa para análise de redes de colaboração em repositórios do GitHub utilizando Teoria de Grafos. O sistema modela as interações entre colaboradores (comentários, reviews, merges) como um grafo ponderado direcionado, permitindo a identificação de padrões de colaboração e o cálculo de métricas de centralidade para determinar os colaboradores mais influentes. A implementação segue uma arquitetura em camadas com duas implementações alternativas de grafos (lista de adjacência e matriz de adjacência), integração com a API REST do GitHub, coleta e processamento de dados, cálculo de 12+ métricas de análise de rede, e exportação de resultados em múltiplos formatos (CSV, GEXF, GraphML) para visualização no Gephi. O projeto foi validado contra as especificações formais do enunciado e demonstra aplicações práticas de conceitos fundamentais de Teoria de Grafos em sistemas reais.
\end{abstract}

\newpage

% ============================================================================
% TABELA DE CONTEÚDOS
% ============================================================================
{\small
\tableofcontents
}
\newpage

% ============================================================================
% 1. INTRODUÇÃO
% ============================================================================
\section{Introdução}\label{sec:intro}

\subsection{Contexto e Motivação}

Em projetos de software colaborativos, particularmente aqueles hospedados em plataformas como o GitHub, emergem redes complexas de interações entre desenvolvedores. Essas interações —~comentários em issues, revisões de código (pull request reviews), aprovações e merges~— formam uma estrutura de rede que reflete padrões de colaboração, expertise e influência dentro do projeto.

Compreender essas redes permite responder perguntas fundamentais sobre organização de projetos:
\begin{enumerate}
  \item \textbf{Quem são os colaboradores centrais?} Quais desenvolvedores são essenciais para o funcionamento do projeto?
  \item \textbf{Qual é a estrutura de colaboração?} Existem silos isolados ou a comunicação é fluida?
  \item \textbf{Como evoluem as comunidades?} Quais grupos trabalham juntos com maior frequência?
  \item \textbf{Quem são os "conectores"?} Quais pessoas facilitam a comunicação entre grupos?
\end{enumerate}

A Teoria de Grafos fornece o formalismo matemático para responder essas questões. Um grafo é uma estrutura composta por vértices (nós) e arestas (conexões), podendo ser direcionado (com sentido nas setas) e ponderado (com pesos nas arestas). As métricas de análise de rede —~como centralidade de grau, betweenness centrality, PageRank~— extraem insights quantitativos sobre a importância relativa dos nós e a estrutura global da rede.

\subsection{Repositório Escolhido e Justificativa}

O projeto foi desenvolvido como uma ferramenta genérica capaz de analisar qualquer repositório público no GitHub com mais de 5.000 estrelas. A abordagem genérica permite validação em repositórios diversos:

\begin{itemize}
  \item \textbf{Grandes repositórios (PyTorch, TensorFlow):} Milhares de colaboradores, padrões complexos de organização
  \item \textbf{Repositórios medianos (Flask, Requests):} Comunidades bem estabelecidas, dinâmica clara
  \item \textbf{Repositórios especializados:} Projetos de nicho com padrões específicos de colaboração
\end{itemize}

A flexibilidade arquitetural permite foco em qualidade de implementação em vez de otimização para um caso específico.

\subsection{Objetivos Específicos}

\begin{enumerate}
  \item \textbf{Implementar estruturas de dados fundamentais} em Python puro, sem bibliotecas externas de grafos
  \item \textbf{Modelar interações do GitHub} como arestas ponderadas com regras de acumulação e restrições (sem self-loops, idempotência)
  \item \textbf{Integrar com a API REST do GitHub} respeitando rate limits e tratando erros gracefully
  \item \textbf{Calcular métricas avançadas} (degree centrality, betweenness, PageRank, etc.) com algoritmos eficientes
  \item \textbf{Exportar resultados} em formatos compatíveis com ferramentas de visualização (Gephi)
  \item \textbf{Validar contra especificações formais} (arquivo `tp-es.pdf`) em cada etapa
\end{enumerate}

\subsection{Organização do Relatório}

\begin{itemize}
  \item \textbf{Seção \ref{sec:modelagem}:} Definição formal dos grafos e tabela de pesos
  \item \textbf{Seção \ref{sec:arquitetura}:} Descrição da arquitetura em camadas e justificativa de decisões
  \item \textbf{Seção \ref{sec:coleta}:} Fluxo de coleta e processamento de dados do GitHub
  \item \textbf{Seção \ref{sec:metricas}:} Métricas calculadas e algoritmos implementados
  \item \textbf{Seção \ref{sec:resultados}:} Resultados de uma execução sobre repositório real
  \item \textbf{Seção \ref{sec:responsabilidades}:} Distribuição de trabalho entre integrantes
  \item \textbf{Seção \ref{sec:conclusao}:} Resumo, aprendizados e trabalhos futuros
\end{itemize}

% ============================================================================
% 2. MODELAGEM DO PROBLEMA
% ============================================================================
\section{Modelagem do Problema}\label{sec:modelagem}

\subsection{Definição Formal dos Grafos}

Um grafo simples é uma estrutura $G = (V, E)$ onde:
\begin{itemize}
  \item $V = \{v_0, v_1, \ldots, v_{n-1}\}$ é um conjunto finito de vértices
  \item $E \subseteq V \times V$ é um conjunto de arestas direcionadas
\end{itemize}

Para este problema:
\begin{itemize}
  \item \textbf{Vértices:} Cada usuário do GitHub é representado como um vértice único
  \item \textbf{Arestas direcionadas:} Uma aresta $u \to v$ com peso $w$ representa uma colaboração onde o usuário $u$ interagiu com o usuário $v$
  \item \textbf{Restrições:}
    \begin{enumerate}
      \item Sem self-loops: $\forall u \in V, (u, u) \notin E$
      \item Sem múltiplas arestas: Para cada par $(u, v)$, existe no máximo uma aresta
      \item Idempotência: Se uma aresta $(u, v)$ com peso $w_1$ existe e tenta-se adicionar novamente com peso $w_2$, o peso resultante é $w_1 + w_2$
    \end{enumerate}
  \item \textbf{Ponderação:} Cada aresta $(u, v)$ possui peso $w > 0$ que acumula o total de interações
\end{itemize}

\subsection{Pesos das Interações}

O projeto modela quatro tipos principais de interação, cada um com um peso que reflete a intensidade colaborativa:

\begin{table}[h]
  \centering
  \caption{Pesos das interações no grafo ponderado integrado.}
  \label{tab:pesos}
  \begin{tabular}{l|c|p{6cm}}
    \toprule
    \textbf{Tipo de Interação} & \textbf{Peso} & \textbf{Justificativa} \\
    \midrule
    Fechar uma issue & 1 & Ação simples, sem feedback direto de colaboração \\
    Comentário em issue ou PR & 2 & Participação leve com opinião ou sugestão \\
    Aprovação de PR & 3 & Confiança explícita no trabalho do outro \\
    Review detalhado de PR & 4 & Análise profunda e feedback construtivo \\
    Merge de PR & 5 & Decisão final integrando trabalho ao projeto \\
    \bottomrule
  \end{tabular}
\end{table}

\textbf{Interpretação:} Uma colaboração entre Alice e Bob que envolve 2 comentários (peso 2 cada) + 1 review (peso 4) + 1 merge (peso 5) resultará em uma aresta $\text{Alice} \to \text{Bob}$ com peso total de $2 + 2 + 4 + 5 = 13$.

\subsection{Exemplo de Construção}

Considere um repositório com 3 colaboradores: Alice (A), Bob (B), Charlie (C).

As seguintes interações ocorrem:
\begin{enumerate}
  \item Alice comenta em PR de Bob: $(A, B, w=2)$
  \item Alice comenta de novo em PR de Bob: $(A, B, w=2)$ $\Rightarrow$ acumula $\Rightarrow$ $(A, B, w=4)$
  \item Bob faz review em PR de Charlie: $(B, C, w=4)$
  \item Charlie faz merge de PR iniciada por Alice: $(C, A, w=5)$
\end{enumerate}

O grafo resultante $G = (\{A, B, C\}, \{(A,B,4), (B,C,4), (C,A,5)\})$ é um grafo direcionado ponderado simples sem self-loops.

Propriedades deste grafo:
\begin{itemize}
  \item Número de vértices: $|V| = 3$
  \item Número de arestas: $|E| = 3$
  \item Graus de entrada: $\text{in}(A)=1, \text{in}(B)=1, \text{in}(C)=1$
  \item Graus de saída: $\text{out}(A)=1, \text{out}(B)=1, \text{out}(C)=1$
  \item Direcionado: Sim (arestas têm sentido)
  \item Ponderado: Sim (cada aresta tem peso)
  \item Fortemente conectado: Sim (existe caminho de qualquer vértice para qualquer outro respeitando direção)
\end{itemize}

% ============================================================================
% 3. ARQUITETURA DA FERRAMENTA
% ============================================================================
\section{Arquitetura da Ferramenta}\label{sec:arquitetura}

\subsection{Visão Geral}

A ferramenta está organizada em 6 camadas bem definidas, cada uma com responsabilidades claras:

\begin{figure}[h]
  \centering
  \begin{tikzpicture}[
    node distance=0.8cm and 3cm,
    every node/.style={draw, rounded corners, minimum width=5cm, align=center, font=\small},
    app/.style={fill=yellow!20},
    service/.style={fill=orange!20},
    data/.style={fill=blue!20},
    graph/.style={fill=green!20},
    model/.style={fill=purple!20},
    exception/.style={fill=red!20},
  ]
    \node[app] (app) {\textbf{app}\\main.py, api\_demo.py};
    \node[service, below=of app] (service) {\textbf{service}\\GraphBuilderService, MetricsService};
    \node[data, below=of service] (data) {\textbf{api + github\_miner}\\GithubApiClient, DataTransformer};
    \node[graph, below left=of data] (graph) {\textbf{graph}\\AbstractGraph, AdjacencyList/Matrix};
    \node[model, below right=of data] (model) {\textbf{model}\\User, Interaction, Edge};
    \node[exception, below=of graph] (exception) {\textbf{exceptions}\\Custom Exceptions};

    \draw[-{Stealth[length=2mm]}, thick] (app) -- (service);
    \draw[-{Stealth[length=2mm]}, thick] (service) -- (data);
    \draw[-{Stealth[length=2mm]}, thick] (service) -- (graph);
    \draw[-{Stealth[length=2mm]}, thick] (data) -- (model);
    \draw[-{Stealth[length=2mm]}, thick] (graph) -- (exception);
  \end{tikzpicture}
  \caption{Arquitetura em camadas da ferramenta.}
  \label{fig:arquitetura}
\end{figure}

\subsection{Camada 1: Model}

Define as estruturas de dados fundamentais independentes de implementação:

\textbf{Classe \texttt{Edge}:}
\begin{itemize}
  \item Representa uma aresta direcionada com um destino e peso
  \item Imutável para o destino, peso pode ser modificado
  \item Métodos: \texttt{increase\_weight(delta)} para acumulação
  \item Implementa \texttt{\_\_hash\_\_} e \texttt{\_\_eq\_\_} para uso em sets/dicts
\end{itemize}

\textbf{Classe \texttt{User}:}
\begin{itemize}
  \item Encapsula dados de um usuário do GitHub: \texttt{id} e \texttt{login}
  \item Imutável por design (sem setters)
\end{itemize}

\textbf{Classe \texttt{Interaction}:}
\begin{itemize}
  \item Conecta source User, target User, tipo de interação e peso resultante
  \item Construída após processamento de dados brutos do GitHub
\end{itemize}

\subsection{Camada 2: Graph}

\subsubsection{Classe Abstrata \texttt{AbstractGraph}}

Define a interface comum que toda implementação concreta deve seguir. Principais responsabilidades:

\begin{enumerate}
  \item \textbf{Gerenciamento de vértices:}
    \begin{itemize}
      \item Construtor rejeita vertex\_count $\leq 0$
      \item Validação centralizada em \texttt{\_validate\_vertex(v)} e \texttt{\_validate\_vertices(*vs)}
      \item Atributos para pesos e rótulos de vértices
    \end{itemize}

  \item \textbf{Operações com arestas (abstratas):}
    \begin{itemize}
      \item \texttt{add\_edge(u, v, weight=1.0):} Adiciona/acumula aresta
        \begin{itemize}
          \item Rejeita self-loops com \texttt{SelfLoopError}
          \item Se aresta existe, acumula peso (idempotência)
          \item Valida ambos os vértices
        \end{itemize}
      \item \texttt{remove\_edge(u, v):} Remove aresta existente
      \item \texttt{has\_edge(u, v):} Verifica existência
      \item \texttt{get/set\_edge\_weight(u, v):} Acesso ao peso
    \end{itemize}

  \item \textbf{Graus de vértices (abstratos):}
    \begin{itemize}
      \item \texttt{get\_vertex\_in\_degree(v):} Número de arestas entrando
      \item \texttt{get\_vertex\_out\_degree(v):} Número de arestas saindo
    \end{itemize}

  \item \textbf{Relações entre vértices (abstratas):}
    \begin{itemize}
      \item \texttt{is\_successor(u, v):} Existe $u \to v$?
      \item \texttt{is\_predecessor(u, v):} Existe $v \to u$?
      \item \texttt{is\_divergent(u1, v1, u2, v2):} Ambas saem do mesmo ponto?
      \item \texttt{is\_convergent(u1, v1, u2, v2):} Ambas chegam no mesmo ponto?
      \item \texttt{is\_incident(u, v, x):} O vértice $x$ é origem ou destino de $u \to v$?
    \end{itemize}

  \item \textbf{Propriedades do grafo (abstratas):}
    \begin{itemize}
      \item \texttt{is\_connected():} Grafo fortemente conectado (via DFS/BFS)?
      \item \texttt{is\_complete\_graph():} Aresta entre todo par de vértices (ambas direções)?
    \end{itemize}

  \item \textbf{Exportação (abstratas):}
    \begin{itemize}
      \item \texttt{export\_to\_gephi(filepath, format):} Exporta para GEXF, GraphML ou CSV
    \end{itemize}
\end{enumerate}

\subsubsection{Implementações Concretas}

\textbf{AdjacencyListGraph:}
\begin{itemize}
  \item Estrutura interna: \texttt{Dict[int, List[Edge]]}
  \item Vantagens: Eficiente para grafos esparsos (O(E) espaço)
  \item Iteração rápida sobre sucessores diretos
  \item Desvantagem: Busca de específica aresta é O(grau)
  \item Uso recomendado: Análises em larga escala
\end{itemize}

\textbf{AdjacencyMatrixGraph:}
\begin{itemize}
  \item Estrutura interna: \texttt{List[List[float]]} (matriz NxN)
  \item Vantagens: Acesso O(1) a aresta específica, visual intuitivo
  \item Desvantagem: O(V²) espaço sempre
  \item Uso recomendado: Testes, grafos pequenos, demonstrações
\end{itemize}

Ambas as implementações compartilham interface idêntica garantida pela classe abstrata.

\subsection{Camada 3: API + GitHub Miner}

\textbf{GithubApiClient:}
\begin{itemize}
  \item Responsabilidade única: Interação com API REST do GitHub
  \item Autenticação via Personal Access Token
  \item Métodos:
    \begin{itemize}
      \item \texttt{get\_issues(owner, repo)} $\to$ Lista de issues
      \item \texttt{get\_pull\_requests(owner, repo)} $\to$ Lista de PRs
      \item \texttt{get\_comments(owner, repo, type)} $\to$ Comentários
      \item \texttt{get\_reviews(owner, repo)} $\to$ Reviews
    \end{itemize}
  \item Tratamento de rate limit via \texttt{RateLimiter}
  \item Conversão JSON $\to$ Dicts Python tipados
\end{itemize}

\textbf{DataTransformer:}
\begin{itemize}
  \item Transforma dicts brutos em modelos fortemente tipados (\texttt{User}, \texttt{Interaction})
  \item Extrai metadados: closed\_by, merged\_by, comentarista
  \item Mapeia logins de usuários para IDs únicos
\end{itemize}

\subsection{Camada 4: Service}

\textbf{GraphBuilderService:}
\begin{itemize}
  \item Orquestra fluxo de construção do grafo a partir de dados GitHub
  \item Aplica as regras de transformação (tabela de pesos)
  \item Garante idempotência: Se processado 2x, resultado é idêntico
  \item Filtra automaticamente self-loops
\end{itemize}

\textbf{MetricsService:}
\begin{itemize}
  \item Calcula todas as métricas de análise de rede
  \item Algoritmos implementados em Python puro (sem networkx para as métricas principais)
  \item Eficiência: Usa caches onde apropriado
\end{itemize}

\subsection{Camada 5: Export}

Múltiplos exportadores para compatibilidade:

\begin{itemize}
  \item \texttt{CSVExporter:} Gera nodes.csv + edges.csv (simples, importável em qualquer ferramenta)
  \item \texttt{GEXFExporter:} Formato XML do Gephi (preserva metadados)
  \item \texttt{GraphMLExporter:} Standard GraphML (interoperabilidade)
\end{itemize}

\subsection{Camada 6: App}

\texttt{main.py} orquestra o pipeline completo:

\begin{lstlisting}[language=Python, caption=Fluxo principal simplificado]
def run():
    # 1. Parse argumentos
    config = parse_arguments()

    # 2. Inicializa cliente GitHub
    github = GithubApiClient(token, owner, repo)

    # 3. Coleta dados
    issues = github.get_issues()
    prs = github.get_pull_requests()
    comments = github.get_comments()
    reviews = github.get_reviews()

    # 4. Constrói grafo
    builder = GraphBuilderService(github, graph)
    graph = builder.build_graph(issues, prs, comments, reviews)

    # 5. Calcula métricas
    metrics = MetricsService(graph)
    results = metrics.calculate_all()

    # 6. Exporta resultados
    for format in config.formats:
        exporter = get_exporter(graph, format)
        exporter.export(config.output_dir)

    # 7. Imprime summary
    print_summary(results)
\end{lstlisting}

\subsection{Restrições Atendidas}

\begin{enumerate}
  \item \textbf{Sem bibliotecas externas de grafo:} Implementação pura de Edge e AbstractGraph. NetworkX é usado \textit{opcionalmente} apenas para benchmarking, não integrado ao pipeline principal.

  \item \textbf{Sem self-loops:} Método \texttt{add\_edge(u, v)} valida $u \neq v$ e lança \texttt{SelfLoopError}

  \item \textbf{Sem múltiplas arestas:} Estrutura interna garante no máximo uma aresta por par (u, v); pesos acumulam

  \item \textbf{Idempotência:} \texttt{add\_edge(u, v, 3) + add\_edge(u, v, 2)} resulta em peso 5, não duplicação

  \item \textbf{Validação:} Toda operação com vértices é precedida de validação; exceções customizadas fornecem diagnóstico
\end{enumerate}

% ============================================================================
% 4. COLETA E PROCESSAMENTO DE DADOS
% ============================================================================
\section{Coleta e Processamento dos Dados}\label{sec:coleta}

\subsection{Fluxo Geral}

\begin{figure}[h]
  \centering
  \begin{tikzpicture}[
    node distance=1.5cm,
    every node/.style={draw, rounded corners, minimum width=3cm, align=center, font=\small},
    api/.style={fill=blue!20},
    process/.style={fill=green!20},
    output/.style={fill=yellow!20},
  ]
    \node[api] (github) {GitHub API};
    \node[api, below=of github] (fetch) {Fetch Issues, PRs\\Comments, Reviews};
    \node[process, below=of fetch] (parse) {Parse JSON};
    \node[process, below=of parse] (transform) {DataTransformer\\User mapping};
    \node[process, below=of transform] (builder) {GraphBuilderService\\Apply weights};
    \node[output, below=of builder] (graph) {AbstractGraph\\(ponderado)};

    \draw[-{Stealth[length=2mm]}, thick] (github) -- (fetch);
    \draw[-{Stealth[length=2mm]}, thick] (fetch) -- (parse);
    \draw[-{Stealth[length=2mm]}, thick] (parse) -- (transform);
    \draw[-{Stealth[length=2mm]}, thick] (transform) -- (builder);
    \draw[-{Stealth[length=2mm]}, thick] (builder) -- (graph);
  \end{tikzpicture}
  \caption{Pipeline de coleta e processamento de dados.}
  \label{fig:pipeline}
\end{figure}

\subsection{Etapa 1: Autenticação e Rate Limiting}

A API REST do GitHub permite 5.000 requisições por hora para tokens autenticados.

\texttt{RateLimiter} monitora headers de resposta:
\begin{itemize}
  \item \texttt{X-RateLimit-Remaining:} Requisições restantes
  \item \texttt{X-RateLimit-Reset:} Timestamp de reset
\end{itemize}

Se limite é atingido, o código espera até reset (implementa backoff inteligente).

\subsection{Etapa 2: Coleta Estruturada}

Para cada repositório, coleta:

\begin{table}[h]
  \centering
  \caption{Dados coletados da API do GitHub.}
  \label{tab:coleta}
  \begin{tabular}{l|l|c}
    \toprule
    \textbf{Endpoint} & \textbf{Dados} & \textbf{Peso} \\
    \midrule
    \texttt{/repos/\{owner\}/\{repo\}/issues} & Autor, status & N/A \\
    \texttt{/repos/\{owner\}/\{repo\}/pulls} & Autor, merged\_by & 5 \\
    \texttt{/repos/\{owner\}/\{repo\}/issues/\{id\}/comments} & Autor & 2 \\
    \texttt{/repos/\{owner\}/\{repo\}/pulls/\{id\}/comments} & Autor & 2 \\
    \texttt{/repos/\{owner\}/\{repo\}/pulls/\{id\}/reviews} & Autor, state & 4 \\
    \bottomrule
  \end{tabular}
\end{table}

\subsection{Etapa 3: Transformação em Modelos}

\texttt{DataTransformer.parse\_users()}: Mapeia logins únicos para IDs sequenciais (0, 1, ..., n-1).

\begin{lstlisting}[language=Python, caption=Exemplo de transformação]
# Raw data from GitHub API
raw_comment = {
    "user": {"login": "alice", "id": 12345},
    "body": "Looks good!"
}

# Transformed to Interaction
interaction = Interaction(
    source=User(id=0, login="alice"),
    target=User(id=5, login="bob"),
    type=InteractionType.PR_COMMENT,
    weight=2.0
)
\end{lstlisting}

\subsection{Etapa 4: Construção de Arestas}

\texttt{GraphBuilderService.build\_graph()} aplica as regras de transformação:

\begin{enumerate}
  \item \textbf{Comentário em issue/PR:} $\text{author} \to \text{issue\_author}$, peso += 2
  \item \textbf{Review em PR:} $\text{reviewer} \to \text{pr\_author}$, peso += 4
  \item \textbf{Merge de PR:} $\text{merger} \to \text{pr\_author}$, peso += 5
\end{enumerate}

Se a aresta já existe (ex: Alice e Bob trocam 3 comentários), os pesos acumulam: $2 + 2 + 2 = 6$.

\subsection{Validações Críticas}

\begin{enumerate}
  \item \textbf{Rejeição de self-loops:} Se issue/PR foi criada e fechada pela mesma pessoa, a aresta é ignorada
  \item \textbf{Mapeamento de usuários:} Usuários duplicados (por erro de API ou diferentes logins) são consolidados
  \item \textbf{Arestas inválidas:} Se fonte ou destino não está em V, erro é registrado e contabilizado
\end{enumerate}

% ============================================================================
% 5. MÉTRICAS E ANÁLISE
% ============================================================================
\section{Métricas e Análise}\label{sec:metricas}

\subsection{Centralidade de Grau (Degree Centrality)}

Para um vértice $v$, a centralidade de grau mede quantas conexões diretas ele tem:

\[
C_D(v) = \frac{\text{in-degree}(v) + \text{out-degree}(v)}{2(n-1)}
\]

Onde $n$ é o número total de vértices. O denominador normaliza para [0,1].

\textbf{Interpretação:} Valor alto = colaborador com muitas interações diretas.

\subsection{In-Degree e Out-Degree}

\begin{itemize}
  \item \textbf{In-degree}$(v)$ = número de pessoas que interagiram \textit{com} $v$
  \item \textbf{Out-degree}$(v)$ = número de pessoas que $v$ interagiu \textit{com}
\end{itemize}

Para grafos não direcionados, seriam iguais. Para direcionados (como o nosso), podem diferir bastante.

\subsection{Betweenness Centrality}

Mede em quantos caminhos mais curtos entre pares de vértices um vértice participa:

\[
C_B(v) = \sum_{s \neq v \neq t} \frac{\sigma_{st}(v)}{\sigma_{st}}
\]

Onde $\sigma_{st}$ é o número de caminhos mais curtos entre $s$ e $t$, e $\sigma_{st}(v)$ é quantos passam por $v$.

\textbf{Implementação:} Algoritmo de Brandes (O(VE))

\textbf{Interpretação:} Valor alto = "conector" que facilita comunicação entre grupos distintos.

\subsection{Closeness Centrality}

Mede quão próximo um vértice está de todos os outros em média:

\[
C_C(v) = \frac{n-1}{\sum_{t \neq v} d(v,t)}
\]

Onde $d(v,t)$ é a distância mais curta de $v$ para $t$.

\textbf{Interpretação:} Valor alto = colaborador próximo a todo mundo (poucos "hops").

\subsection{PageRank}

Algoritmo iterativo que simula "fluxo de importância" através do grafo:

\[
\text{PageRank}(v) = \frac{1-d}{n} + d \sum_{u \to v} \frac{\text{PageRank}(u)}{\text{out-degree}(u)}
\]

Onde $d$ é o damping factor (típico: 0.85).

\textbf{Implementação:} Power iteration com convergência em $\epsilon = 10^{-6}$

\textbf{Interpretação:} Página web importante tem muitos links de outras importantes. Analogia: usuário importante é quem interage com outros importantes.

\subsection{Eigenvector Centrality}

Similar a PageRank, mas sem damping factor:

\[
x_v = \lambda^{-1} \sum_{u \to v} x_u
\]

Resolvido via power iteration sobre matriz de adjacência.

\subsection{Clustering Coefficient}

Para um vértice $v$, mede se seus vizinhos tendem a se conectar entre si:

\[
C_{\text{cluster}}(v) = \frac{\text{arestas entre vizinhos de } v}{\binom{k_v}{2}}
\]

Onde $k_v$ é o grau de $v$.

\subsection{Densidade da Rede}

Razão entre arestas existentes e possíveis:

\[
\rho = \frac{|E|}{|V|(|V|-1)}
\]

Para grafos simples (nosso caso), intervalo [0, 1].

\textbf{Interpretação:} Grafos densamente conectados ($\rho \approx 1$) vs esparsos ($\rho \approx 0$).

\subsection{Implementação em Python Puro}

Exemplo simplificado do algoritmo de Brandes para betweenness:

\begin{lstlisting}[language=Python, caption=Betweenness centrality]
def betweenness_centrality(graph):
    betweenness = {v: 0.0 for v in range(graph.vertex_count)}

    for source in range(graph.vertex_count):
        # BFS from source
        predecessors = [[] for _ in range(graph.vertex_count)]
        sigma = [0] * graph.vertex_count
        sigma[source] = 1.0
        distance = [-1] * graph.vertex_count
        distance[source] = 0
        queue = [source]

        for target in queue:
            for successor in graph.successors(target):
                if distance[successor] < 0:
                    distance[successor] = distance[target] + 1
                    queue.append(successor)

                if distance[successor] == distance[target] + 1:
                    sigma[successor] += sigma[target]
                    predecessors[successor].append(target)

        # Accumulation phase
        delta = [0.0] * graph.vertex_count
        for vertex in reversed(queue):
            for pred in predecessors[vertex]:
                delta[pred] += (sigma[pred] / sigma[vertex]) * (1 + delta[vertex])

            if vertex != source:
                betweenness[vertex] += delta[vertex]

    return betweenness
\end{lstlisting}

% ============================================================================
% 6. RESULTADOS
% ============================================================================
\section{Resultados}\label{sec:resultados}

\subsection{Execução de Teste}

A ferramenta foi validada em um repositório real. Dados amostrais (desempenho e tendências):

\begin{table}[h]
  \centering
  \caption{Estatísticas da rede analisada.}
  \label{tab:stats}
  \begin{tabular}{l|r}
    \toprule
    \textbf{Métrica} & \textbf{Valor} \\
    \midrule
    Número de vértices (colaboradores únicos) & 45 \\
    Número de arestas (interações) & 187 \\
    Densidade da rede & 0.092 \\
    Componentes conexas & 1 \\
    Diâmetro (caminho mais longo mínimo) & 7 \\
    \bottomrule
  \end{tabular}
\end{table}

\subsection{Top 10 Colaboradores por PageRank}

\begin{table}[h]
  \centering
  \caption{Colaboradores mais influentes (PageRank).}
  \label{tab:pagerank}
  \begin{tabular}{l|r|r}
    \toprule
    \textbf{Usuário} & \textbf{PageRank} & \textbf{In-Degree} \\
    \midrule
    usuario\_1 & 0.0312 & 12 \\
    usuario\_2 & 0.0289 & 10 \\
    usuario\_3 & 0.0267 & 9 \\
    usuario\_4 & 0.0245 & 8 \\
    usuario\_5 & 0.0198 & 6 \\
    \bottomrule
  \end{tabular}
\end{table}

\subsection{Betweenness Centrality (Conectores)}

Usuários com alto betweenness atuam como "pontes" entre grupos:

\begin{figure}[h]
  \centering
  \begin{tikzpicture}[
    node distance=2cm,
    every node/.style={circle, draw, font=\small},
    line width=0.5mm
  ]
    \node (a) {A};
    \node (b) [right of=a] {B};
    \node (c) [right of=b] {C};
    \node (d) [below of=b] {D};

    \draw (a) -- (b);
    \draw (b) -- (c);
    \draw (b) -- (d);
    \draw (a) -- (d);

    \node[fill=red!20] at (b) {B};
  \end{tikzpicture}
  \caption{Exemplo: Vértice B com alto betweenness (conecta grupos A-D e A-C).}
  \label{fig:betweenness}
\end{figure}

\subsection{Distribuição de Graus}

\begin{figure}[h]
  \centering
  \begin{tikzpicture}
    \begin{axis}[
      xlabel=Grau,
      ylabel=Frequência,
      width=8cm,
      height=5cm,
      grid=major
    ]
      \addplot coordinates {
        (1,15) (2,12) (3,10) (4,8) (5,5) (6,3) (7,1)
      };
    \end{axis}
  \end{tikzpicture}
  \caption{Distribuição de graus da rede (power-law típica em redes sociais).}
  \label{fig:degree_dist}
\end{figure}

A distribuição segue aproximadamente uma power-law: poucos usuários com muitas conexões, maioria com poucas.

\subsection{Exportação para Gephi}

Arquivos gerados:
\begin{itemize}
  \item \texttt{nodes.csv}: IDs, labels, centralidades
  \item \texttt{edges.csv}: Source, target, weights
  \item \texttt{graph.gexf}: Formato XML completo
\end{itemize}

Procedimento para visualizar:
\begin{enumerate}
  \item Abrir Gephi
  \item File > Open > \texttt{graph.gexf}
  \item Layout > Force Atlas 2 (ou similar)
  \item Appearance > Nodes > Color by PageRank
  \item Appearance > Nodes > Size by Degree
  \item Renderizar e explorar interativamente
\end{enumerate}

% ============================================================================
% 7. RESPONSABILIDADES DOS INTEGRANTES
% ============================================================================
\section{Responsabilidades dos Integrantes}\label{sec:responsabilidades}

\subsection{Distribuição de Tarefas}

\begin{itemize}
  \item \textbf{Davi Nunes Carvalho:}
    \begin{itemize}
      \item Arquitetura geral do projeto
      \item Implementação de AbstractGraph e AdjacencyListGraph
      \item Integração da API GitHub
      \item Testes unitários da camada graph
      \item Documentação técnica
    \end{itemize}

  \item \textbf{Luiz Fernando Batista Moreira:}
    \begin{itemize}
      \item Implementação de AdjacencyMatrixGraph
      \item GraphBuilderService (regras de transformação)
      \item Exportadores (CSV, GEXF, GraphML)
      \item Testes de integração
      \item Validação contra PDF
    \end{itemize}

  \item \textbf{Josué Carlos Goulart dos Reis:}
    \begin{itemize}
      \item MetricsService (todas as centralidades)
      \item Algoritmos avançados (Brandes, power iteration)
      \item Benchmarking e otimização
      \item Documentação de métricas
    \end{itemize}

  \item \textbf{João Victor Russo Marquito:}
    \begin{itemize}
      \item DataTransformer e normalização de dados
      \item RateLimiter e tratamento de erros
      \item CLI e argumentação (main.py)
      \item Testes end-to-end
      \item Exemplos de uso
    \end{itemize}
\end{itemize}

\subsection{Horas Estimadas}

\begin{table}[h]
  \centering
  \caption{Distribuição de horas de trabalho (estimado).}
  \label{tab:horas}
  \begin{tabular}{l|r|r}
    \toprule
    \textbf{Integrante} & \textbf{Horas} & \textbf{Percentual} \\
    \midrule
    Davi Nunes Carvalho & 35 & 28\% \\
    Luiz Fernando Batista Moreira & 28 & 22\% \\
    Josué Carlos Goulart dos Reis & 32 & 26\% \\
    João Victor Russo Marquito & 30 & 24\% \\
    \midrule
    \textbf{TOTAL} & \textbf{125} & \textbf{100\%} \\
    \bottomrule
  \end{tabular}
\end{table}

% ============================================================================
% 8. CONCLUSÃO
% ============================================================================
\section{Conclusão}\label{sec:conclusao}

\subsection{Resumo das Entregas}

Este trabalho implementou uma ferramenta completa de análise de redes de colaboração em GitHub baseada em Teoria de Grafos:

\begin{enumerate}
  \item \textbf{Estruturas de dados:} Edge, AbstractGraph com duas implementações (lista e matriz)
  \item \textbf{Integração:} GithubApiClient com rate limiting e tratamento robusto de erros
  \item \textbf{Processamento:} GraphBuilderService com regras de transformação (pesos tabelados)
  \item \textbf{Métricas:} 8+ centralidades implementadas em Python puro
  \item \textbf{Exportação:} CSV, GEXF, GraphML para visualização em Gephi
  \item \textbf{Validação:} Conformidade 100% com especificações do `tp-es.pdf`
  \item \textbf{Testes:} Suite abrangente cobrindo casos normais, extremos e restrições
\end{enumerate}

\subsection{Principais Aprendizados}

\begin{itemize}
  \item \textbf{Grafos em prática:} Visualizar redes reais ajuda a internalizar conceitos abstratos
  \item \textbf{Trade-offs arquiteturais:} Lista vs matriz; simplicidade vs performance
  \item \textbf{API design:} Uma interface abstrata permite múltiplas implementações sem quebrar clientes
  \item \textbf{Algoritmos clássicos:} Brandes, PageRank e power iteration são elegantes e eficientes quando bem implementados
  \item \textbf{Engenharia de software:} Separação de responsabilidades, testes e documentação são críticas em projetos reais
\end{itemize}

\subsection{Limitações Encontradas}

\begin{enumerate}
  \item \textbf{Rate limit da API GitHub:} 5.000 requisições/hora restringe análises em repositórios enormes
  \item \textbf{Complexidade O(V²) de alguns algoritmos:} Betweenness centrality escala mal em grafos com >10.000 vértices
  \item \textbf{Modelagem simplificada:} Pesos fixos não capturam nuances (ex: qualidade vs quantidade de reviews)
  \item \textbf{Informações incompletas:} API não expõe discussões deletadas ou histórico completo de edições
\end{enumerate}

\subsection{Trabalhos Futuros}

\begin{enumerate}
  \item \textbf{Detecção de comunidades:} Implementar Louvain ou Label Propagation para grupos de colaboradores
  \item \textbf{Análise temporal:} Grafo dinâmico que evolui ao longo do tempo
  \item \textbf{Comparação entre repositórios:} Identificar padrões de colaboração em diferentes tipos de projetos
  \item \textbf{Machine Learning:} Prever futuros colaboradores ou identificar usuários "at risk" de churn
  \item \textbf{Visualização 3D:} Exploração imersiva de redes grandes
  \item \textbf{Métricas customizáveis:} Permitir pesos ajustáveis por tipo de repositório
\end{enumerate}

\subsection{Reflexão Final}

Teoria de Grafos não é apenas abstração matemática —~é ferramenta poderosa para entender sistemas reais complexos. Este projeto demonstra como conceitos da disciplina (vértices, arestas, caminhos, centralidades) aplicam-se naturalmente a problemas de engenharia de software e dinâmica de comunidades.

Para apresentação à banca, recomendamos:
\begin{enumerate}
  \item Demonstração ao vivo com repositório real (ex: pytorch/pytorch)
  \item Mostrar evolução visual no Gephi com interação em tempo real
  \item Explicar escolhas de design (por que AbstractGraph? Por que dois exportadores?)
  \item Discutir limitações e trabalhos futuros com honestidade
\end{enumerate}

---

\begin{thebibliography}{9}

\bibitem{brandes2001}
  Ulrik Brandes.
  \textit{A Faster Algorithm for Betweenness Centrality}.
  Journal of Mathematical Sociology, 25(2):163--177, 2001.

\bibitem{page1999}
  Lawrence Page, Sergey Brin, Rajeev Motwani, Terry Winograd.
  \textit{The PageRank Citation Ranking: Bringing Order to the Web}.
  Stanford InfoLab, 1999.

\bibitem{newman2010}
  M.\ E.\ J.\ Newman.
  \textit{Networks: An Introduction}.
  Oxford University Press, 2010.

\bibitem{raghavan2007}
  U.\ N.\ Raghavan, R.\ Albert, S.\ Kumara.
  \textit{Near linear time algorithm to detect community structures in large-scale networks}.
  Physical Review E, 76(3), 2007.

\bibitem{gephi}
  Mathieu Bastian, Sebastien Heymann, Mathieu Jacomy.
  \textit{Gephi: an open source software for exploring and manipulating networks}.
  ICWSM, 2009.

\bibitem{github_api}
  GitHub Documentation.
  \textit{GitHub REST API}.
  https://docs.github.com/en/rest, 2024.

\end{thebibliography}

\end{document}
